\begin{lemma}[The one-step drift]
Let $F$ be a finite-dimensional real inner-product space, let
$1\le C_0<\infty$ and $0\le\kappa<\infty$, and let $v\in\Aclass_{m,C_0}$ in the
sense of \noderef{IsFitClass} be centered at $(0,1)$ with constant $C_0$ in the
sense of \noderef{IsCenteredAt} and satisfy the horizontal Jensen bound with
constant $\kappa$ at every time in the sense of \noderef{HasJensenBound}. Write
$f=W$ for the contact envelope of \noderef{ContactEnv}, $f^\ell=\Penv_\ell f$
for the forward envelope of \noderef{QuadEnvP}, $g_\epsilon=\Menv_\epsilon v$
for the backward envelope of \noderef{QuadEnvM}, and $J_\epsilon$ for the
quantity of \noderef{JGap}. Let $\epsilon>0$, let $z=(x,t)$ with
$t>\epsilon$, and let $(z^+,z^-)$ be an optimizing step pair of size $\epsilon$
at $z$ in the sense of \noderef{IsStepPair}, with $z^\pm=(x^\pm,t\pm\epsilon)$.
Put
\[
  Q_+=\frac{x^+-x}{\epsilon},\qquad
  Q_-=\frac{x-x^-}{\epsilon},
\]
and let $Q_0\in F$ be the reference velocity of
\noderef{DPSpatialRegularity} at $z$ and length $\epsilon$, i.e.\ a vector at
which $f^\epsilon(\cdot,t)$ has gradient $Q_0$ at $x$:
\[
  Q_0=\nabla_xf^\epsilon(x,t).
\]
Define
\begin{align*}
  D&=f(z)+f^{2\epsilon}(z)-2f^\epsilon(z),\\
  e&=f^\epsilon(z^-)+\frac{\epsilon}{2}|Q_-|^2-f(z),\\
  \widetilde D&=2g_\epsilon(z)-v(z)-\Menv_{2\epsilon}v(z),\\
  \widetilde e&=v(z)+\frac{\epsilon}{2}|Q_+|^2-g_\epsilon(z^+).
\end{align*}
Then all four quantities are nonnegative, and
\begin{equation}\label{eq:dp-drift}
\begin{split}
  \frac{J_\epsilon(z^+)+J_\epsilon(z^-)}{2}-J_\epsilon(z)
  &=\frac{D+e+\widetilde D+\widetilde e}{4\epsilon}\\
  &\ge\frac{D+e}{4\epsilon}\\
  &\ge c\bigl(|Q_+-Q_0|^2+|Q_--Q_0|^2\bigr),
  \qquad c=\frac{1}{16(1+\log2)}.
\end{split}
\end{equation}
\end{lemma}
\begin{proof}
Throughout, $z=(x,t)$ with $t>\epsilon>0$, $x^\pm$ are the spatial parts of
$z^\pm$, and $f=W$ denotes the contact envelope of \noderef{ContactEnv}, so
that $f(y,\tau)=\Cenv_\tau v(y,\tau)-\kappa^2\tau/2$ for $\tau>0$. Note
$t-\epsilon>0$ and $t+\epsilon>0$, so every value of $f$ written below is taken
at a positive time, and $v$ is defined at every time.

\emph{Step 0: standing facts about $v$ and $f$.}
By \noderef{IsFitClass} the function $v$ is continuous on all of space-time and
is a viscosity solution of $F[v]=0$ there; in particular it is a viscosity
supersolution on all of space-time in the sense of \noderef{IsLimitSuperOn} and
a subsolution in the sense of \noderef{IsLimitSubOn}, by
\noderef{IsLimitSolutionOn}. By \noderef{UniformSubquadratic}, applied with the
centering point $(0,1)$, $v$ has subquadratic spatial growth on compact time
intervals of $\R$ in the sense of \noderef{SubquadraticGrowth}; the same
property with $I=(0,\infty)$ follows, since shrinking $I$ only shrinks the
family of compact subintervals tested. By \noderef{BaryEnvSandwich}, applied
with $I=\R$, $f=v$ and the constant $\kappa$, we have for every $\sigma>0$ and
every $(y,\tau)$
\begin{equation}\label{eq:dpd-sandwich}
  0\le\Cenv_\sigma v(y,\tau)-v(y,\tau)\le\frac{\kappa^2\sigma}{2},
\end{equation}
and taking $\sigma=\tau>0$ with \noderef{ContactEnv} gives
\begin{equation}\label{eq:dpd-fsandwich}
  v(y,\tau)-\frac{\kappa^2\tau}{2}\le f(y,\tau)\le v(y,\tau)\qquad(\tau>0).
\end{equation}
Consequently $f$ has subquadratic spatial growth on compact subintervals of
$(0,\infty)$: given $J\subseteq(0,\infty)$ compact and $\varepsilon'>0$, put
$\tau_{\max}=\max J$, apply \noderef{SubquadraticGrowth} for $v$ with $J$ and
$\varepsilon'/2$ to obtain $M'>0$, and enlarge $M'$ so that
$\kappa^2\tau_{\max}/2\le(\varepsilon'/2)M'^2$; then
\eqref{eq:dpd-fsandwich} gives $|f(y,\tau)|\le\varepsilon'|y|^2$ for
$|y|\ge M'$ and $\tau\in J$. Moreover $f$ is continuous on
$\{(y,\tau):\tau>0\}$: by \noderef{BaryEnvDuality}, part (3), applied to $v$
with $I=\R$, the map $(\sigma,\zeta)\mapsto\Cenv_\sigma v(\zeta)$ is continuous
at every pair with $\sigma>0$; composing with the continuous map
$\zeta\mapsto(\tau,\zeta)$, where $\tau$ is the time coordinate of $\zeta$, and
subtracting the continuous function $\kappa^2\tau/2$ gives the claim. Finally
$v$ is a viscosity subsolution on the open set $\{\tau>0\}$: if $\phi$ is twice
continuously differentiable at a point $\zeta$ of positive time and $v-\phi$ has
a local maximum at $\zeta$ relative to $\{\tau>0\}$, then, $\{\tau>0\}$ being
open and containing $\zeta$, a neighbourhood witnessing that relative maximum
may be intersected with $\{\tau>0\}$ and is then a neighbourhood of $\zeta$ in
the whole space on which the inequality holds, so $v-\phi$ has a local maximum
at $\zeta$ in the whole space and $0\le F[\phi](\zeta)$ follows. Hence
\noderef{BarycentricSubsolution}, applied to $v$ with this subsolution
property, with the subquadratic growth on $(0,\infty)$ and with
$c=\kappa^2/2$, shows that
\begin{equation}\label{eq:dpd-fsub}
  f\text{ is a viscosity subsolution of }F[\cdot]=0\text{ on }\{\tau>0\},
\end{equation}
since $\Cenv_\tau v(\zeta)-(\kappa^2/2)\tau$ is exactly $f$ by
\noderef{ContactEnv}.

\emph{Claim B: the envelopes are genuine suprema and infima.}
Let $I\subseteq\R$, let $g$ be continuous on $\{(y,\tau):\tau\in I\}$ with
subquadratic spatial growth on compact subintervals of $I$ in the sense of
\noderef{SubquadraticGrowth}, let $\ell>0$ and let $\zeta=(\xi,\tau)$ with
$\tau\in I$.
\begin{enumerate}
\item[(i)] If $\tau+\ell\in I$, the set
$\{g(y,\tau+\ell)-|y-\xi|^2/(2\ell):y\in F\}$ of \noderef{QuadEnvP} is nonempty
and bounded above, so $\Penv_\ell g(\zeta)$ is finite, is at least every member
of that set, and is at most every upper bound for it.
\item[(ii)] If $\tau-\ell\in I$, the set
$\{g(y,\tau-\ell)+|y-\xi|^2/(2\ell):y\in F\}$ of \noderef{QuadEnvM} is nonempty
and bounded below, so $\Menv_\ell g(\zeta)$ is finite, is at most every member,
and is at least every lower bound.
\end{enumerate}
\emph{Proof.} (i) The choice $y=\xi$ shows the set is nonempty. Apply
\noderef{SubquadraticGrowth} with the compact set $J=\{\tau+\ell\}\subseteq I$
and $\varepsilon'=1/(4\ell)>0$ to get $M>0$ with
$|g(y,\tau+\ell)|\le|y|^2/(4\ell)$ for $|y|\ge M$; and $g(\cdot,\tau+\ell)$ is
continuous, hence bounded in absolute value by some $C_M$ on the closed ball of
radius $M$. From $|y|^2\le2|y-\xi|^2+2|\xi|^2$ we get
$|y-\xi|^2\ge|y|^2/2-|\xi|^2$, so for every $y$
\[
  g(y,\tau+\ell)-\frac{|y-\xi|^2}{2\ell}
   \le C_M+\frac{|y|^2}{4\ell}-\frac{|y|^2}{4\ell}+\frac{|\xi|^2}{2\ell}
   =C_M+\frac{|\xi|^2}{2\ell},
\]
an upper bound independent of $y$. (ii) Apply (i) to $-g$, which is continuous
with the same growth property, at the time $\tau-\ell$, and multiply the
resulting bound by $-1$. This proves Claim B.

By Step 0, Claim B applies to $g=f$ with $I=(0,\infty)$ at every center of
positive time whose shifted time is also positive, and to $g=v$ with $I=\R$ at
every center and length.

\emph{Step 1: $D\ge0$.}
Apply \noderef{QuadEnvPathConvex}, part (1), to $f$ with $I=(0,\infty)$,
$T=2\epsilon>0$, the subsolution property \eqref{eq:dpd-fsub}, the continuity
and growth of Step 0, the base point $x_0=x$, the velocity $d'=0$ and the time
$t$; its hypothesis $t+\ell\in(0,\infty)$ for $\ell\in[0,2\epsilon]$ holds since
$t>0$. Since $\|d'\|=0$, the conclusion is that
$\ell\mapsto\Penv_\ell f(x,t)$ is convex on $[0,2\epsilon]$. The midpoint
inequality at $\ell=0$ and $\ell=2\epsilon$, whose midpoint is $\epsilon$,
gives
\[
  \Penv_\epsilon f(x,t)
   \le\tfrac12\Penv_0f(x,t)+\tfrac12\Penv_{2\epsilon}f(x,t),
\]
and $\Penv_0f(x,t)=f(z)$ by \noderef{QuadEnvP}. Multiplying by $2$ and
rearranging gives $f(z)+f^{2\epsilon}(z)-2f^\epsilon(z)\ge0$, i.e.\ $D\ge0$.

\emph{Step 2: $\widetilde D\ge0$.}
Apply \noderef{QuadEnvPathConvex}, part (2), to $v$ with $I=\R$,
$T=2\epsilon>0$, the supersolution property of Step 0, and the continuity and
growth of Step 0, at the point $(x,t)$; its hypothesis $t-\ell\in\R$ is
vacuous. The conclusion is that $\ell\mapsto\Menv_\ell v(x,t)$ is concave on
$[0,2\epsilon]$, so the midpoint inequality gives
\[
  \Menv_\epsilon v(x,t)
   \ge\tfrac12\Menv_0v(x,t)+\tfrac12\Menv_{2\epsilon}v(x,t),
\]
and $\Menv_0v(x,t)=v(z)$ by \noderef{QuadEnvM}. Multiplying by $2$ and
rearranging gives $2g_\epsilon(z)-v(z)-\Menv_{2\epsilon}v(z)\ge0$, i.e.\
$\widetilde D\ge0$.

\emph{Step 3: $e\ge0$, and a competitor bound used later.}
The center $z^-=(x^-,t-\epsilon)$ has positive time and
$(t-\epsilon)+\epsilon=t>0$, so Claim B(i) applies to $f$ at $z^-$ with length
$\epsilon$. For every $y\in F$ the point $y$ is a member of the set of
\noderef{QuadEnvP} for $f^\epsilon(z^-)$, whose time argument is $t$; hence
\begin{equation}\label{eq:dpd-competitor}
  f^\epsilon(z^-)\ge f(y,t)-\frac{|y-x^-|^2}{2\epsilon}\qquad(y\in F).
\end{equation}
Taking $y=x$ and using $x-x^-=\epsilon Q_-$, so that
$|x-x^-|^2/(2\epsilon)=\epsilon|Q_-|^2/2$, gives
$f^\epsilon(z^-)\ge f(z)-\epsilon|Q_-|^2/2$, i.e.\ $e\ge0$.

\emph{Step 4: $\widetilde e\ge0$.}
The center $z^+=(x^+,t+\epsilon)$ has $(t+\epsilon)-\epsilon=t$, and Claim
B(ii) applies to $v$ at $z^+$ with length $\epsilon$ and $I=\R$. The point
$y=x$ is a member of the set of \noderef{QuadEnvM} for $g_\epsilon(z^+)$, whose
time argument is $t$, so
\[
  g_\epsilon(z^+)\le v(z)+\frac{|x-x^+|^2}{2\epsilon}
   =v(z)+\frac{\epsilon|Q_+|^2}{2},
\]
using $x^+-x=\epsilon Q_+$. This is exactly $\widetilde e\ge0$.

\emph{Step 5: the drift identity.}
By \noderef{DPStepIdentities}, applied with the present data, the two half-step
identities hold:
\begin{equation}\label{eq:dpd-half}
  f^{2\epsilon}(z)=f^\epsilon(z^+)-\frac{\epsilon}{2}|Q_+|^2,\qquad
  \Menv_{2\epsilon}v(z)=g_\epsilon(z^-)+\frac{\epsilon}{2}|Q_-|^2 .
\end{equation}
The first gives
$f^\epsilon(z^+)=f^{2\epsilon}(z)+\epsilon|Q_+|^2/2$, and the definition of
$e$ gives $f^\epsilon(z^-)=f(z)+e-\epsilon|Q_-|^2/2$. Adding and subtracting
$2f^\epsilon(z)$,
\begin{equation}\label{eq:dpd-fsum}
  f^\epsilon(z^+)+f^\epsilon(z^-)-2f^\epsilon(z)
   =D+e+\frac{\epsilon}{2}\bigl(|Q_+|^2-|Q_-|^2\bigr).
\end{equation}
Likewise the second identity in \eqref{eq:dpd-half} gives
$g_\epsilon(z^-)=\Menv_{2\epsilon}v(z)-\epsilon|Q_-|^2/2$, and the definition
of $\widetilde e$ gives
$g_\epsilon(z^+)=v(z)+\epsilon|Q_+|^2/2-\widetilde e$. Adding and subtracting
$2g_\epsilon(z)$,
\begin{equation}\label{eq:dpd-gsum}
  g_\epsilon(z^+)+g_\epsilon(z^-)-2g_\epsilon(z)
   =-\widetilde D-\widetilde e
     +\frac{\epsilon}{2}\bigl(|Q_+|^2-|Q_-|^2\bigr),
\end{equation}
because $-\widetilde D=v(z)+\Menv_{2\epsilon}v(z)-2g_\epsilon(z)$. The half-step
cost difference $\epsilon(|Q_+|^2-|Q_-|^2)/2$ is the same in
\eqref{eq:dpd-fsum} and \eqref{eq:dpd-gsum}, so subtracting the second from the
first cancels it:
\[
  \bigl(f^\epsilon-g_\epsilon\bigr)(z^+)
   +\bigl(f^\epsilon-g_\epsilon\bigr)(z^-)
   -2\bigl(f^\epsilon-g_\epsilon\bigr)(z)
   =D+e+\widetilde D+\widetilde e .
\]
By \noderef{JGap}, $J_\epsilon=(f^\epsilon-g_\epsilon)/(2\epsilon)$, so
dividing by $4\epsilon$ turns the left side into
$\bigl(J_\epsilon(z^+)+J_\epsilon(z^-)\bigr)/2-J_\epsilon(z)$ and proves the
first line of \eqref{eq:dp-drift}. Since $\widetilde D\ge0$ and
$\widetilde e\ge0$ by Steps 2 and 4, and $4\epsilon>0$, the second line
follows.

\emph{Step 6: the velocity estimate.}
Let $b,d\in F$ be arbitrary. Apply \noderef{QuadEnvPathConvex}, part (1), to
$f$ with $I=(0,\infty)$, $T=2\epsilon$, the base point $x_0=x+d$, the velocity
$d'=(b-d)/(2\epsilon)$ and the time $t$; the hypotheses are those of Step 1.
Writing
\[
  \Phi(\ell)=\Penv_\ell f\bigl(x+d+\ell d',t\bigr)
     +\frac{|b-d|^2}{4\epsilon^2}\,\ell\log\ell
   \qquad(0\le\ell\le2\epsilon),
\]
where $\|d'\|^2=|b-d|^2/(4\epsilon^2)$ and $\ell\log\ell$ takes the value $0$
at $\ell=0$, the conclusion is that $\Phi$ is convex on $[0,2\epsilon]$. Now
$x+d+0\cdot d'=x+d$, $x+d+2\epsilon d'=x+b$ and
$x+d+\epsilon d'=x+(b+d)/2$, and $\Penv_0f=f$ by \noderef{QuadEnvP}, so
\[
  \Phi(0)=f(x+d,t),\qquad
  \Phi(2\epsilon)=f^{2\epsilon}(x+b,t)+\frac{|b-d|^2}{2\epsilon}\log(2\epsilon),
\]
\[
  \Phi(\epsilon)=f^\epsilon\Bigl(x+\frac{b+d}{2},t\Bigr)
     +\frac{|b-d|^2}{4\epsilon}\log\epsilon .
\]
The midpoint inequality $2\Phi(\epsilon)\le\Phi(0)+\Phi(2\epsilon)$ and
$\log(2\epsilon)-\log\epsilon=\log2$ therefore give
\begin{equation}\label{eq:dp-path-midpoint}
  f(x+d,t)+f^{2\epsilon}(x+b,t)
    -2f^\epsilon\Bigl(x+\frac{b+d}{2},t\Bigr)
  \ge-\frac{\log2}{2\epsilon}|b-d|^2 .
\end{equation}

We bound the left side of \eqref{eq:dp-path-midpoint} from above by three
inequalities. First, \eqref{eq:dpd-competitor} with $y=x+d$, together with
$x+d-x^-=d+\epsilon Q_-$ and
\[
  \frac{|d+\epsilon Q_-|^2}{2\epsilon}
   =\frac{|d|^2}{2\epsilon}+Q_-\cdot d+\frac{\epsilon}{2}|Q_-|^2,
\]
and with $f^\epsilon(z^-)=f(z)+e-\epsilon|Q_-|^2/2$ from the definition of $e$,
gives
\begin{equation}\label{eq:dpd-first}
  f(x+d,t)\le f(z)+e+Q_-\cdot d+\frac{|d|^2}{2\epsilon}.
\end{equation}
Second, \noderef{DPSpatialRegularity}, applied at the time $t>0$, the length
$2\epsilon>0$ and the center $x$, provides the unique maximizer $y_{2\epsilon}$
with $\nabla_xf^{2\epsilon}(x,t)=(2\epsilon)^{-1}(y_{2\epsilon}-x)$ and the two
supporting inequalities at that gradient. By \noderef{DPStepIdentities} the
same function has gradient $Q_+$ at $x$, and a function has at most one
gradient at a point, so $(2\epsilon)^{-1}(y_{2\epsilon}-x)=Q_+$. The upper
supporting inequality with $h=b$ is therefore
\begin{equation}\label{eq:dpd-second}
  f^{2\epsilon}(x+b,t)\le f^{2\epsilon}(z)+Q_+\cdot b+\frac{|b|^2}{2t}.
\end{equation}
Third, \noderef{DPSpatialRegularity} at the time $t$, the length
$\epsilon>0$ and the center $x$ provides the unique maximizer $y_\epsilon$ with
$\nabla_xf^\epsilon(x,t)=\epsilon^{-1}(y_\epsilon-x)$ and the same two
inequalities; by hypothesis $f^\epsilon(\cdot,t)$ has gradient $Q_0$ at $x$, so
again by uniqueness of the gradient $\epsilon^{-1}(y_\epsilon-x)=Q_0$. The
lower supporting inequality with $h=(b+d)/2$ is
\[
  f^\epsilon(z)+Q_0\cdot\frac{b+d}{2}-\frac{|b+d|^2}{8\epsilon}
   \le f^\epsilon\Bigl(x+\frac{b+d}{2},t\Bigr),
\]
that is, after multiplying by $-2$,
\begin{equation}\label{eq:dpd-third}
  -2f^\epsilon\Bigl(x+\frac{b+d}{2},t\Bigr)
   \le-2f^\epsilon(z)-Q_0\cdot(b+d)+\frac{|b+d|^2}{4\epsilon}.
\end{equation}
Adding \eqref{eq:dpd-first}, \eqref{eq:dpd-second} and \eqref{eq:dpd-third},
and using $D=f(z)+f^{2\epsilon}(z)-2f^\epsilon(z)$, the left side of
\eqref{eq:dp-path-midpoint} is at most
\[
  D+e+(Q_+-Q_0)\cdot b+(Q_--Q_0)\cdot d
       +\frac{|b|^2}{2t}+\frac{|d|^2}{2\epsilon}
       +\frac{|b+d|^2}{4\epsilon}.
\]
Combining with \eqref{eq:dp-path-midpoint} and rearranging,
\[
  D+e\ge(Q_0-Q_+)\cdot b+(Q_0-Q_-)\cdot d
    -\frac{|b|^2}{2t}-\frac{|d|^2}{2\epsilon}
    -\frac{|b+d|^2}{4\epsilon}-\frac{\log2}{2\epsilon}|b-d|^2 .
\]
Since $t>\epsilon$ we have $|b|^2/(2t)\le|b|^2/(2\epsilon)$, and
$|b\pm d|^2\le2(|b|^2+|d|^2)$ gives
$|b+d|^2/(4\epsilon)\le(|b|^2+|d|^2)/(2\epsilon)$ and
$(\log2/(2\epsilon))|b-d|^2\le(\log2/\epsilon)(|b|^2+|d|^2)$. The four
subtracted terms therefore total at most
\[
  \Bigl(\frac{1}{2\epsilon}+\frac{1}{2\epsilon}+\frac{\log2}{\epsilon}\Bigr)
     \bigl(|b|^2+|d|^2\bigr)
  =\frac{1+\log2}{\epsilon}\bigl(|b|^2+|d|^2\bigr),
\]
so that
\begin{equation}\label{eq:dpd-free}
  D+e\ge(Q_0-Q_+)\cdot b+(Q_0-Q_-)\cdot d
     -\frac{1+\log2}{\epsilon}\bigl(|b|^2+|d|^2\bigr)
\end{equation}
for all $b,d\in F$. Put $A=Q_0-Q_+$, $B=Q_0-Q_-$,
$\mu=\epsilon/(2(1+\log2))>0$, and choose $b=\mu A$, $d=\mu B$ in
\eqref{eq:dpd-free}. Then
\[
  D+e\ge\mu\bigl(|A|^2+|B|^2\bigr)
     -\frac{1+\log2}{\epsilon}\mu^2\bigl(|A|^2+|B|^2\bigr)
   =\Bigl(\mu-\frac{\mu}{2}\Bigr)\bigl(|A|^2+|B|^2\bigr)
   =\frac{\epsilon}{4(1+\log2)}\bigl(|A|^2+|B|^2\bigr),
\]
where the middle equality uses
$((1+\log2)/\epsilon)\mu^2=\mu/2$. Dividing by $4\epsilon>0$ and recalling
$|A|=|Q_+-Q_0|$, $|B|=|Q_--Q_0|$,
\[
  \frac{D+e}{4\epsilon}
   \ge\frac{1}{16(1+\log2)}\bigl(|Q_+-Q_0|^2+|Q_--Q_0|^2\bigr),
\]
which is the last line of \eqref{eq:dp-drift} with $c=1/(16(1+\log2))$. This
completes the proof.
\end{proof}
